% report.tex: the technical report. Every number comes from tables/ (generated by rule tables) and
% every figure from figs/ (rule figures); nothing here is typed by hand. The manuscript is written
% from this document and checked against it (WORKFLOW.md section 5, principle 2).
\documentclass[11pt]{article}
\usepackage[margin=1in]{geometry}
\usepackage{booktabs}
\usepackage{graphicx}
\usepackage{amsmath,amssymb}
\usepackage{caption}
\usepackage{hyperref}
\usepackage{xcolor}
\input{tables/numbers.tex}

\title{Textbook priors over visual features:\\ what a vision-language model's textbook knowledge is worth in labelled images}
\author{Technical report, generated by the workflow}
\date{\today}

\begin{document}
\maketitle

\section{Question and design}

A multimodal language model carries textbook knowledge about what pathology looks like. On the
\nDatasets{} MedMNIST v2 2D benchmarks \cite{yang2023}, a vision-language model (\primaryModel) scored
every image of a seeded test sample against a cited bank of diagnostic visual features (the concept
bank; \texttt{data/concepts/}, built with a simulated expert review and no clinician). Five arms
predict on the same test sample per dataset: A, the model's zero-shot class distribution; B, the
nearest textbook fingerprint over the concept scores, with no labels; C, an L2 logistic regression on
the concept scores with $n$ labels; P, the same regression on frozen ImageNet ResNet-18 features
with the same $n$ labels and the same nested subsets; and E, a ResNet-18 trained from scratch on the
full official split. The primary metric is test AUC from the \texttt{medmnist} evaluator; ACC is
reported but decides nothing. Inference across datasets is a one-sided sign test. The three
hypotheses and their decision rules were fixed before any number existed (\texttt{WORKFLOW.md},
section 2). \callsArchived{} model calls are archived under \texttt{results/score/}. \flaggedCells{}
dataset--model cells exceeded the 5\% incompleteness gate and are marked $\dagger$.

\section{Reconciliation of arm E against the published table}

Before anything else is trusted: arm E, trained with the published recipe, against the ResNet-18
rows of the MedMNIST v2 benchmark, on the full official test split. \reconAgree{} of \reconRows{}
cells (of \reconExpected{} expected) agree within tolerance.

\begin{center}\small\input{tables/reconciliation.tex}\end{center}

\section{H1: substitution}

\emph{Textbook features are worth a measurable number of labelled images.} At $n=\firstN$, arm C
beats arm P on \hOneWins{} of \hOneN{} datasets (one-sided sign test $p=\hOneP$). The median $n_B$,
the labelled images the pixel probe needs to reach the zero-label textbook arm, is
\nBmedian. \textbf{H1 is \hOneVerdict.}

\begin{center}\small\resizebox{\textwidth}{!}{\input{tables/h1.tex}}\end{center}

\begin{figure}[h]\centering
\includegraphics[width=\textwidth]{figs/fig_curve.png}
\caption{The learning curve per dataset: arm C (concept scores) against arm P (ImageNet features)
at equal labels, with arm B's zero-label line and its interval, arm E's ceiling, and the crossing
$n_B$ (marker).}
\end{figure}

\begin{figure}[h]\centering
\includegraphics[width=0.8\textwidth]{figs/fig_nb.png}
\caption{$n_B$ per dataset with its 95\% bootstrap interval, ordered by modality: where the
textbook pays. Arrows mark crossings censored at the grid edge.}
\end{figure}

\section{H2: the bank, not just the model}

\emph{Directing the model at cited visual features beats asking it for the diagnosis, and the bank's
structure carries the difference.} Arm B beats arm A on \hTwoWins{} of \hTwoN{} datasets
($p=\hTwoP$); under the permutation controls, \hTwoControls. \textbf{H2 is \hTwoVerdict.}

\begin{center}\small\resizebox{\textwidth}{!}{\input{tables/h2.tex}}\end{center}

\section{H3: scale}

\emph{The prior gets better with a bigger model.} Arm B's AUC over four models and the arm-B
datasets, analysed as one blocked $2\times2$ design (family $\times$ size tier, dataset as block),
with the planned linear contrast on $\log_{10}$ parameters and a Friedman test with Nemenyi
post-hoc. \textbf{H3 is \hThreeVerdict.}

\begin{center}\small\input{tables/ladder.tex}\end{center}
\begin{center}\small\input{tables/ladder_anova.tex}\end{center}

\begin{figure}[h]\centering
\includegraphics[width=\textwidth]{figs/fig_ladder.png}
\caption{The model ladder. Left: arm B's AUC per dataset as measured, which shows why datasets are
never pooled. Right: the same values centred within dataset, with the across-dataset mean and its
95\% interval; the only panel from which a trend should be read.}
\end{figure}

\section{Completeness of the response archive}

The fraction of sampled images with at least one missing concept answer (or an unusable zero-shot
reply), per dataset and model--split; cells above 5\% are flagged $\dagger$ and excluded from the
headline counts.

\begin{center}\small\resizebox{\textwidth}{!}{\input{tables/completeness.tex}}\end{center}

\section{What is not claimed}

The concept scores are not clinically validated; the simulated review does not substitute for a
clinician; no arm is state of the art; and since every source dataset is public and labelled,
``the model carries textbook knowledge'' and ``the model has seen this benchmark'' cannot be
separated here. The bank's sources follow.

\nocite{*}
\bibliographystyle{plain}
\bibliography{references}

\end{document}
